\chapter{Construction Parameters, Registrations, and Corrections}\label{app:tuning-protocol}

This appendix records three things a reader may need to check: the
hyperparameters the frozen construction ran at, the decision rules that were
fixed in dated documents before the runs they govern, and the figures this
thesis withdraws from earlier versions of its own analysis.

\section{Canonical Hyperparameters}\label{sec:app-hyperparams}

\begin{table}[htbp]
\centering
\caption{Construction hyperparameters, read from the configuration of the
frozen run every result is attributed to (Appendix~\ref{app:frozen-tree}).}
\label{tab:hyperparams}
\begin{tabular}{lll}
\toprule
Parameter & Symbol/Name & Value \\
\midrule
Embedding dimension & $d$ & 256 (fixed) \\
Maximum recursion depth & $D_{\max}$ & 8 \\
Minimum cluster size & $n_{\min}$ (\texttt{minClusterSize}) & 55 \\
Proposal separation bar & $\epsilon_{\text{sep}}$ (\texttt{proposalSeparationBar}) & 0.025 \\
Acceptance tolerance & $\tau$ (\texttt{tau}) & $10^{-6}$ \\
Acceptance gate, in paired SE units & \texttt{acceptanceZ} & 2.0 \\
Membership share floor & \texttt{membershipFloor} & 0.25 \\
Routing beam margin & \texttt{routingBeamGamma} & 0.20 \\
Descent-gate slack & \texttt{descentMargin} & 0.12 \\
Effective support floor & \texttt{effectiveSupportFloor} & 2.0 \\
Max leaf assignments (arena-time cap) & \texttt{maxLeafAssignments} & 5 \\
Fusion pre-filter threshold & \texttt{fusionSimilarityThreshold} & 0.90 \\
Branching cap & \texttt{maxK} & 4 \\
Iteration budget & \texttt{numIterations} & 50 \\
Construction mode & \texttt{dagMode} & \texttt{DAG\_MAX} \\
\bottomrule
\end{tabular}
\end{table}

Four entries are commonly read backwards and are stated here in the direction
the evidence supports.

\texttt{minClusterSize} is a \textbf{budget} choice. It sets the cell count to
what the judge budget can identify; per-cell reliability follows from it rather
than motivating it. It was set directly against the judge-call budget once
granularity turned out to be evaluatively flat
(Section~\ref{sec:res-granularity}), a result that made the choice free; its
value of $55$ was not selected by any sweep.

\texttt{descentMargin} is a \textbf{mode selector}, not a tuned optimum. At
$0.12$ the descent gate holds no query above a leaf, so the frozen artifact has
an exhaustive leaf partition and no reserve of unexplained mass. The boundary
below which the gate begins holding queries has not been located at this
configuration, so neither it nor the cost of crossing it is quoted. What is
measured is that $0.12$ over-provisions: the quantity the margin pays for is
the gap between a node's fitted direction and its children's resultant, whose
worst value over twenty-one runs is $0.065$, so the margin is $1.85\times$ the
largest gap ever seen (Appendix~\ref{app:dag-formalism}).

\texttt{numIterations} is an upper budget, not a run length. The frozen
construction reached its certified fixed point at \textbf{iteration 10}, and
iterations 8--10 are identical on every recorded column
(Section~\ref{sec:res-construction}). The budget never binds.

\texttt{dagMode} carries no polyhierarchy setting; the name is historical and
cross-linking is no longer part of the pipeline
(Section~\ref{sec:poly-negative}).

\paragraph{How the values were chosen, and the leakage boundary.} An
L9($3^4$) orthogonal-array sweep with Pareto selection exists in this project's
tooling, but it fixed none of the values above: it was run on a predecessor
parameterisation whose routing layer has since been replaced, and three of its
four factors no longer exist. The surviving values were set outside it, for the
reasons given above. The machinery is retained for the tuning study of the
current parameter set, which is future work.

What matters for validity is a property of that machinery that does hold: no
reserved-set label, accuracy, or judge verdict is read by any selection gate at
any point. This is the mechanism that enforces the tuning/test boundary in
practice rather than asserting it. Because construction is deterministic at a
fixed seed, comparisons across configurations are exact rather than confounded
by run-to-run variance.

\section{Ablations}\label{sec:appendix-ablations}

Construction is deterministic at a fixed seed, so each ablation is a single run
at seed 42 rather than a mean over seeds. One ablation has a current, citable
outcome: the separation-threshold sweep
$\epsilon_{\text{sep}} \in \{0.01, 0.025, 0.05, 0.10, 0.30\}$, which fixed the
canonical $0.025$ between the two nulls of
Section~\ref{sec:arch-requirements}.

The rest are withdrawn or historical, and their numbers cannot be quoted. The
embedding-dimension row could not have been run at all --- the slice width is
a literal in the source with no configuration binding --- so the claim that
$d = 256$ is a verified stable operating point is withdrawn. The
marginal-separation, branching-cap, and $k$-fallback figures predate a
correction to the splitter's routing check and came from runs differing from
canonical in more than one parameter; the qualitative finding they support ---
that the branching cap was a structural determinant before the fallback
existed and a cost bound afterwards --- stands, but no value from them is
quoted. Ablations of the nat-margin assignment gates, the cross-link operator,
and the dynamic routing temperature concern mechanisms that have since been
removed; the cross-link result is reported as a negative result in
Section~\ref{sec:poly-negative}.

\section{Pre-Registrations}\label{app:preregs}

Five decision rules were fixed in dated documents before the data they govern
existed. This section records what each fixed, so that every use of
``pre-registered'' in Chapter~\ref{ch:results} can be checked. Dates from the
version-control history: P1, P2 and P3 on 2026-07-27; P4 on 2026-07-27, amended
through 2026-07-30, every amendment predating the runs it governs; P2 carries a
results-written-back addendum of 2026-07-28.

The purpose is not ceremony. Three results in this project inverted after the
fact, and only those whose reading had been fixed in advance survived the
inversion as findings rather than as revisions.

\begin{table}[htbp]
\centering
\small
\caption{Registration status by paired domain. ``Replication'': the
run used the identical protocol without a prospective registration of
its own. Pre-settlement outcomes came from the superseded
twelve-model batch, so the outcome column records status under the
addendum (Section~\ref{sec:prereg-addendum-p4}), not per-domain
verdicts.}
\label{tab:domain-registration}
\begin{tabular}{lllp{2.0cm}p{2.4cm}}
\toprule
Domain & Registered & Date & Prediction & Outcome \\
\midrule
mathematics & P4 body & 2026-07-27 & null & rule withdrawn$^{\dagger}$ \\
law & P4 body & 2026-07-27 & advantage & rule withdrawn$^{\dagger}$ \\
philosophy & --- & replication & --- & ran; same pattern \\
history & --- & replication & --- & ran; same pattern \\
psychology & Addendum~2 & 2026-07-29 & advantage & rule withdrawn$^{\dagger}$ \\
engineering & Addendum~2 & 2026-07-29 & null & rule withdrawn$^{\dagger}$ \\
physics & Addendum~3 & 2026-07-30 & advantage & rule withdrawn$^{\dagger}$ \\
computer science & Addendum~3 & 2026-07-30 & null & rule withdrawn$^{\dagger}$ \\
mathematics, re-run & Addendum~4 & 2026-07-30 & null & open$^{\ddagger}$ \\
\addlinespace
chemistry & P4 body & 2026-07-27 & --- & registered, never run;
retired with no claim either way \\
\bottomrule
\end{tabular}
\par\smallskip
\footnotesize $^{\dagger}$The registered decisive rule
($|\Delta\rho| \geq 0.007$) is withdrawn
(Section~\ref{sec:prereg-addendum-p4}), so no per-domain
``met''/``failed'' verdict is issued. What the settled batch measures
in every registered domain, under the statistic the addendum promotes:
both arms produce the same ranking (one adjacent transposition in
computer science), the cell-scoped arm abstains more ($13.0\%$ against
$10.9\%$ pooled, direction $8/8$), and agreement conditional on
committing differs by one point ($86.0\%$ against $87.2\%$). The
directional \emph{advantage} predictions therefore find no support and
the \emph{null} predictions are consistent with what was measured ---
stated as description, not as verdicts under the withdrawn rule.
$^{\ddagger}$The settled batch includes a mathematics run, but under a
protocol the registration predates; the registered prediction is open
rather than answered.
\end{table}

\paragraph{P1 --- Granularity and discriminative power.} Fixed three named
outcomes in advance, so any result the analysis could return already had an
interpretation attached, and the rule that the \emph{observed} between-cell
correlation is primary and the disattenuated one secondary. The outcome that
occurred was named in advance as ``the same signal, thinner''. The primacy rule
proved decisive: the disattenuated column moves in the opposite direction
across granularities and clips past $1.0$, so leading with it would have
produced a conclusion from the arithmetic of the correction rather than from
the data (Section~\ref{sec:res-granularity}).

\paragraph{P2 --- Rubric specificity.} Fixed four directional predictions, each
with the direction of its own confound named separately, together with
within-arm decision thresholds and the rule that if the sharpest discriminator
came back at zero in both arms the test is \emph{inconclusive} rather than
negative. As registered it recorded ``treatment arm only, no null, this is not
yet evidence''; the dated addendum writes the null-arm results back in and
records that the null-arm artifacts are untracked in version control
(Section~\ref{sec:res-rubric-null}).

\paragraph{P3 --- Arena launch cost.} Fixed cost bands for per-cell
identification, with the instruction that no band is renegotiated once the
number is seen, and a void condition: a cost constant across cells would
indicate the scheduler was still capping and would report nothing. The measured
cost varies across cells and its minimum falls inside the band
(Section~\ref{sec:res-cost}).

\paragraph{P4 --- The no-partition baseline.} Fixed the paired design, the
primary outcome (Spearman $\rho$ against ground-truth accuracy per domain), the
secondary (per-verdict answer-key agreement), a decisive difference of two
Spearman grid steps, the designation of mathematics as the null arm and law as
the test, and four void conditions --- identical question sets, the confidence
gate disabled in both arms, identical query samplers, and identically defined
ties across the two verdict formats. Recorded in the same document, before the
run: the honest prior that with the multiple-choice options present the
comparison would likely come back null, because the judge can verify the answer
and the rubric is then decorative.

\emph{Two of P4's provisions did not survive later checking.} They are
withdrawn in the dated addendum below rather than edited into the text above,
because a registration revised after seeing its own results stops being one.

\paragraph{P5 --- The options-blind arm.} A fully options-blind arm has not been
run, but its predictions were fixed in writing per stratum and are reproduced
because a prediction written before a run that never happened is still evidence
of discipline. Predicted: the untraced-against-untraced stratum should return to
chance with a very high tie rate, and a result at $70\%$ or above would indicate
leakage, with the candidates enumerable --- question recognition from
pre-training, position bias surviving the dual-call control, or answer-letter
frequency priors. The alternative was written down too: landing unchanged would
mean the judge was never using the key on traced pairs, which is a cleaner
decomposition than any currently available and should not be discovered after
the fact.

The settled batch runs in \texttt{RESOLVED} mode, which injects the answering
model's selected option as text rather than withholding the options block.
That mode is adjacent to this registration but is not the registered
options-blind arm, and no result in this thesis reports it as satisfying P5.

\section{Register of Corrections}\label{app:corrections-register}

Earlier versions of this pipeline and of its analysis reported figures this
thesis corrects or withdraws. Each is told once, here. \textbf{No value computed
before a correction is quoted anywhere in this thesis.}

\paragraph{The centred-residual quartet.} An earlier draft reported a
single quartet of centred between-cell residuals ($-0.119$, $-0.024$,
$-0.024$, $-0.024$); no one centring convention reproduces all four
values, so the quartet is withdrawn. The roster-specific values that
replace it, from \texttt{tools/analysis/discriminative\_flatness.py},
are $+0.024$ (14 domains) and $-0.024$ (87 leaves) at the 8-model
roster, $-0.109$ and $-0.009$ at the 11-model band, against mechanical
centring nulls of $-0.077$ and $-0.011$; they have not themselves been
independently confirmed (Section~\ref{sec:res-granularity}).

\paragraph{Bradley--Terry standard errors.} The zero-sum rank completion
requires subtracting $1/K^{2}$ after inversion; the implementation subtracted
$1/K$ (Appendix~\ref{app:numerics}); fixed as commit \texttt{a92a3ed},
2026-07-27. The over-subtraction of $(K-1)/K^{2}$ ---
$0.109$ at $K = 8$ --- exceeds the variance itself at most realistic comparison
counts, so no standard error reported before the fix is a measurement. Where the
over-subtraction drove the variance negative the value was clamped at a floor of
$0.01$; where it did not, the value survived but understated. Correct values are
44 to 209 times larger. This invalidates every confidence interval and every
resolved-pair determination computed before the fix, since pair resolution
thresholds on $3(\sigma_i + \sigma_j)$. No pre-correction interval appears in
this thesis, so there is nothing to restate.

\paragraph{The reliability constant.} Published at $c = 7.66$; the corrected
value is $c = 2.52$. Two independent errors compound. The curve was fitted
against raw split-half correlations, which measure reliability at \emph{half}
length, so fitting $n/(n + 2c_0)$ against $n/(n + c)$ returns $c = 2c_0$; and
$c$ is a property of the roster rather than of the corpus, so re-fitting at the
corrected band multiplies it by a further $0.66$. Separately, disattenuation
must be two-sided --- both sides of $\rho(\text{arena}, \text{ground truth})$
are finite-sample rankings --- so the correction is $\rho / r$, not
$\rho / \sqrt{r}$.

The consequences invert rather than shift. At a $40.8$-query cell an observed
median $\rho$ of $0.823$ corrects to $0.978$ under the published constant and to
$0.874$ under the corrected one; the reliability ladder across granularities
compresses from a spread of $0.219$ to $0.088$, and leaf-scale cells stop being
individually low-reliability. That weakens the reliability leg of the argument
for judging at a coarser cut, leaving the identification-cost leg of
Section~\ref{sec:res-cost} to carry it. Two caveats travel with any value of
$c$: the interval originally quoted for it resamples three fitted points without
propagating sampling uncertainty in the underlying estimates, so it is not a
confidence interval; and the curve was fitted on noise-free ground-truth
rankings, so arena reliability at a given $n$ lies \emph{below} it. It is a
ceiling, not a prediction. It has also not been refitted to the arena roster
(Appendix~\ref{app:models}).

\paragraph{The first evaluative-diversity permutation null.} The first version
sampled pseudo-leaves independently, giving them 11--26\% pairwise question
overlap where real leaves have 1.6--4.9\%. Overlapping pseudo-leaves agree with
each other for free, so the null was inflated and the real cells looked
distinctive against it. The corrected version matches pseudo-leaf overlap to the
real distribution and returns the opposite answer: no domain shows sub-structure
after correction for multiple testing, and three domains' leaves agree
\emph{more} than random splits of the same sizes
(Section~\ref{sec:res-leaf-null}). The first version's conclusion is withdrawn.
It is the cleanest instance in this project of a null whose comparison arm was
never itself checked.

\paragraph{Four figures withdrawn without replacement.} Each was reported by an
earlier analysis, none reproduces, and no claim in this thesis rests on any of
them. \emph{Evaluative-redundancy percentages} of the form ``$72.2\%$ of leaf
pairs redundant'' could not be reproduced under either reading, and the
convention that produced them is underspecified in the pre-registration.
\emph{The correctness-blind rank correlation}, originally a fall to
$0.74$--$0.76$, recomputes to $0.90$--$0.93$ under both tie conventions; the
magnitude is withdrawn and the direction carries the claim
(Section~\ref{sec:res-correctness}). \emph{A leaf-level between-cell agreement
of $0.922$ on re-routed assignments} does not reproduce; the $0.922$ in
Table~\ref{tab:granularity} is a different quantity on construction
assignments, sharing three digits and nothing else. \emph{Routing ECE},
exported at $0.2114$, is not the metric its definition specifies because of an
aggregation defect (Appendix~\ref{app:metrics}); it is quoted only as a
diagnostic and has not been re-exported.

\paragraph{A side-channel found before it could enter a run.} Recorded because
it belongs to the same audit trail, although no reported figure was affected. A
pre-launch format check on two candidate roster entries showed exactly one
anomalous column: newline count, median $0$ where every other model showed 8 to
24. Their stored output was not prose but a raw JSON envelope carrying a
self-reported \texttt{reason\_code} field alongside the reasoning, and that
field predicts the outcome being judged --- $83.2\%$ correct at code A against
$36.5\%$ at code C for one system, $90.0\%$ against $57.5\%$ for the other. The
rendering path returns stored output verbatim, so a judge would have received
the code inline. Neither system entered any arena run as stored.

\section{Addendum to P4, 2026-07-31}\label{sec:prereg-addendum-p4}

Two provisions of P4 are withdrawn, with the measurement that forced each.

\paragraph{1. The generic arm does not vary the partition.} P4 says both arms
judge an identical question set ``because here the partition is the treatment''.
It is not. The generic arm is a replay: it re-judges the exact
$(\text{cell}, \text{pair}, \text{question})$ triples the cell-scoped arm
scheduled, under the same cells. And the two aggregation paths P4 treats as
different --- rolling per-cell fits up by pooling sufficient statistics, against
fitting once at the domain --- are the same computation, because pooling sums
wins and ties into a key that carries no cell identity. Refitting both ways on
identical verdicts gives identical orderings and a maximum $|\theta|$ difference
of $0.000 \times 10^{0}$.

Both arms therefore share the partition, the schedule and the aggregation, and
differ only in the judge prompt. What P4 registered is a
\textbf{rubric-specificity} test at matched budget. No condition in this thesis
removes the partition; a test that did would have to schedule its own
comparisons over the whole domain with no cells, holding the rubric fixed, and
was not run.

\paragraph{2. The decisive difference is smaller than the noise it was meant to
sit above.} P4 fixes $|\Delta\rho| \geq 0.007$, two steps of the Spearman grid
at twelve models. Two runs of an identical configuration --- same domain, same
taxonomy, same question pool, same option mode, differing only in the scheduler
seed --- produced $\Delta\rho$ of $+6.00$ and $0.00$ grid steps. A threshold at
two steps cannot adjudicate a quantity whose run-to-run spread is three to four
times larger.

$\rho$ is therefore reported for continuity and is not used as a decision rule.
The primary outcome becomes the secondary one P4 already names --- per-verdict
answer-key agreement --- which has the power the rank correlation lacks and
which reproduced across the two seeds to within $0.1$ percentage points.

\paragraph{What is not withdrawn.} The paired design, the four void conditions,
the directional predictions and their dates all stand. The registered
prediction for mathematics stands as recorded; whether it holds is open rather
than answered, because the run that answered it is superseded. Recording a
prediction that fails, and a rule that turns out to be unusable, is the point of
registering them.

\section{The Addenda Sequence and Withdrawn Design Baselines}
\label{sec:app-addenda-record}

Demoted from Section~\ref{sec:exp-prereg}, which keeps the summary of
what each registration fixed.

\paragraph{The five addenda, in sequence.} The first is the
results-written-back addendum to P2 (2026-07-28), described with P2 in
Section~\ref{app:preregs}. Of the four that extend P4, each but the
last was written before the runs it governs: the second (2026-07-29)
registers psychology as a reordering domain and engineering as a null.
The third (2026-07-30) registers physics as the last untested flagged
domain and computer science as a third null, and refits the screen at
the then-current twelve-model roster --- a refit itself since
superseded by the settled eight-model one. The fourth (2026-07-30)
registers the mathematics re-run against the coverage confound below,
and fixes in advance how each of its three possible outcomes is to be
read. The fifth (2026-07-30) is written after seven runs have results;
it records the outcomes against the predictions, retires chemistry as
registered but never run, and states explicitly that philosophy and
history were never registered prospectively and are reported as
replications.

\paragraph{The mathematics coverage confound.} In the withdrawn
twelve-model campaign, mathematics was the only run predating the two
scheduler repairs --- the mandatory bootstrap coverage floor with its
in-run assertion, and per-(leaf, pair) query
shuffling;\footnote{Commits \texttt{5e1be09} and \texttt{a5c36fd}
respectively.} its comparison graph was a leader-centric star covering
under a third of the model pairs, and no property of mathematics could
be inferred from it. That campaign is withdrawn in its entirety, so
the confound is recorded here as the reason the re-run registration
exists rather than as a qualification on any reported number. The
fourth addendum registers a re-run on the repaired scheduler under the
original null prediction and fixes three readings in advance: a
reproduced negative means the failed null is real and now rests on a
complete graph; a null means the original was a scheduler artifact and
is superseded on stated grounds; a positive means the screen failed in
the opposite direction from philosophy. Reporting only the re-run is
excluded under every outcome. That run was outstanding when
Chapter~\ref{ch:experimental-design} was fixed.

\paragraph{Design baselines that are not reported.} The baseline
conditions for taxonomy quality specified in an earlier design ---
flat $k$-means, agglomerative clustering, and a random taxonomy null,
scored by a family of reference-free hierarchy metrics --- are not
reported. The supervised hierarchical-classification baselines were
removed as unimplemented; no gold taxonomy exists for MMLU-Pro, so
reference-dependent metrics were dropped rather than computed against
an ad-hoc structure; and the remaining runs were never executed. None
of those metrics bears on any of the three links of
Section~\ref{sec:research-questions}, and the construction evidence
reported instead is the fixed-point certificate, the leaf-size
distribution, and the separation nulls of
Section~\ref{sec:arch-requirements}.

\paragraph{A correction to Addendum 3.} The registration for physics
and computer science described them as having 12 and 3 leaves, with
pools of 546 and 208 questions. Those figures come from the
label-composition count that Chapter~\ref{ch:methodology} sets aside:
they assign each leaf to its majority MMLU-Pro category rather than to
its anchor. Under the adapted partition the runs actually used, the
figures are 10 and 6 cells over pools of 338 and 191. The error is
descriptive and touches no prediction --- both registrations turn on
the sign and magnitude of $\Delta\rho$, neither of which references a
cell count. The secondary observation registered for physics also
survives: mathematics still has the most cells, physics still has the
second most, and the gap is one cell under either count. The
correction is recorded here rather than repaired in the registration,
which is left as written.

\paragraph{Which decision band governs which quantity.} The two
pre-registration documents state bands for different things, and the
results-written-back addendum resolves the assignment. The
$\geq 0.05$-with-$\geq 75\%$-of-cells-positive band governs the
treatment arm's within-arm specificity, met at $0.138$ with $93\%$ of
cells positive. The $\leq 0.05$ / $\geq 0.10$ band governs the random
arm's absolute own-cell specificity, met at $0.012$, and is the band
Section~\ref{sec:res-rubric-null} quotes.

\paragraph{Housekeeping on the rubric-specificity document.} The
results-written-back addendum of 2026-07-28 (described with P2 in
Section~\ref{app:preregs}) closes the document: the headline
separation ($p = 1.21 \times 10^{-6}$) was reproduced from the stored
artifacts on 2026-07-27, and the null-arm artifacts sit in an
untracked build directory --- a reproducibility caveat the document
carries explicitly (Section~\ref{sec:res-rubric-null}).
